
\usepackage{graphicx}  % Please import figures that are *high resolution* PDFs
\usepackage{caption}
\usepackage{subfigure}  % Useful for subfigures
\usepackage{booktabs}  % Useful for high quality tables (e.g., you can replace \hrule with \toprule, \midrule, and \bottomrule).
\usepackage{multicol}
\usepackage{multirow}
%%%%%%%%%%%%%%%%%%%%%%%%%%%%%%%%%%%%%%%%%%%%%%%%%%%%%%%%%%%%%%%%%%%%%%%%%%%%%%%
% MATH PACKAGES (Comment out this section if unnecessary for your dissertation)
%
\usepackage{amsfonts}
\usepackage{amsmath}
\usepackage{amssymb}
\usepackage{amstext}
\usepackage{amsthm}

\usepackage{xcolor}         % colors
\usepackage[capitalize]{cleveref}
%\usepackage{tikz}
\usepackage{amsmath}
\usepackage{breqn}
\usepackage{bbm}
\usepackage{amsthm}
\usepackage{amssymb}
\usepackage{mathtools}
\usepackage{xspace}
\usepackage{enumitem}
\usepackage[hang,flushmargin]{footmisc}
%\usepackage{algpseudocode}
\usepackage{wrapfig}
\setlist{leftmargin=*}
\setlist[itemize]{noitemsep}
\setlist[enumerate]{noitemsep}

\usepackage{listings}
\usepackage{framed}



\newcommand{\predSeq}{\mathbf{s}}
%===============================
\newcommand{\methodname}{CSL\xspace}
\newcommand{\methodnamePrompt}{\texttt{CSL}\xspace}
\newcommand{\methodnameNext}{\texttt{CSL-Next}\xspace}
% CSL : Contextualized Sequence Likelihood

%\newcommand{\methodSampled}{Ours(Sampled)\xspace}
%\newcommand{\methodMLE}{Ours(MostLikely)\xspace}

\newcommand{\baselinePTrue}{\texttt{P(true)}\xspace}
\newcommand{\baselineSAR}{\texttt{TokenSAR}\xspace}
\newcommand{\baselineNLLNorm}{\texttt{SL}(norm)\xspace}
\newcommand{\baselineNLL}{\texttt{SL}\xspace}

\newcommand{\baselineSE}{\texttt{SE}\xspace}
\newcommand{\baselineSENorm}{\texttt{SE(norm)}\xspace}
\newcommand{\methodSEAttn}{\texttt{SE+CSL}\xspace}

\newcommand{\datasetcoqa}{\texttt{coqa}\xspace}
\newcommand{\datasetnqopen}{\texttt{nq}\xspace}
\newcommand{\datasettrivia}{\texttt{trivia}\xspace}

%\newcommand{\baselineEcc}{\texttt{Ecc}\xspace}
\newcommand{\baselineDegree}{\texttt{Deg}\xspace}

\newcommand{\ConfMethod}{\text{C}_{CSL}\xspace}
\newcommand{\Conf}{\text{C}\xspace}
\newcommand{\Acc}{acc\xspace}

\newcommand{\LLaMATwoName}{LLaMA2-70B\xspace}
\newcommand{\GPTName}{\texttt{gpt-3.5-turbo-0125}\xspace}




%\definecolor{falured}{rgb}{0.7, 0.15, 0.15}
%\newcommand\covfail[1]{{\color{falured}#1}}

%==========
\newcommand{\defeq}{\vcentcolon=}
\DeclarePairedDelimiter\ceil{\lceil}{\rceil}
\DeclarePairedDelimiter\floor{\lfloor}{\rfloor}
\newcommand\fontred[1]{{\color{red}#1}}
\newcommand\reminder[1]{{\color{green}#1}}
\DeclareMathOperator*{\argmax}{arg\,max}
\DeclareMathOperator*{\argmin}{arg\,min}

\newcommand{\todo}{\fontred{TODO:}}


\newcommand\attnplus[1]{{\color{red}#1}}
%\newcommand\attnplus[1]{\textbf{#1}}


\newtheorem{theorem}{Theorem}
\newtheorem{proposition}{Proposition}

%\usepackage{floatrow}%floatrow causes caption to be on bottom
%\floatsetup[table]{capposition=top}


\definecolor{codegreen}{rgb}{0,0.6,0}
\definecolor{codegray}{rgb}{0.5,0.5,0.5}
\definecolor{codepurple}{rgb}{0.58,0,0.82}
\definecolor{backcolour}{rgb}{0.95,0.95,0.92}
\lstdefinestyle{mystyle}{
    backgroundcolor=\color{backcolour},   
    commentstyle=\color{codegreen},
    keywordstyle=\color{magenta},
    numberstyle=\tiny\color{codegray},
    stringstyle=\color{codepurple},
    basicstyle=\ttfamily\footnotesize,
    breakatwhitespace=false,         
    breaklines=true,                 
    captionpos=b,                    
    keepspaces=true,                 
    % numbers=left,                    
    numbersep=5pt,                  
    showspaces=false,                
    showstringspaces=false,
    showtabs=false,                  
    tabsize=2,
commentstyle=\color{blue},
morecomment=[l]{$},
}
\lstset{style=mystyle}


